
Nous introduisons le concept de $target\_tiles$(tuile objective). \`A présent pour rentrer dans la tuile d'arrivé et gagner la partie, on a besoin des ressources \begin{scriptsize}NUM\_RESOURCES-2, NUM\_RESOURCES-1 et NUM\_RESOURCES\end{scriptsize}; pour  \begin{scriptsize}NUM\_RESOURCES=4\end{scriptsize} cela correspond a 1 or, 1 gem et 1 clé. Ces ressources sont réparties dans le manoir sur les $target\_tiles$.

Chaque $target\_tile$ est associée à une unique "target ressource". Ainsi, pour gagner, il faut passer par tous les $target\_tiles$.

\subsubsection{Changement en game\_state}

Pour garder en mémoire les $target\_tiles$, on ajoute dans \texttt{game\_state} un tableau de localisation représentant les $target\_tiles$: \texttt{target\_locations}. Chaque indice du tableau est associée une $target\_ressource$: \texttt{target\_locations[0]} est la position où est situé l'or, \texttt{target\_locations[1]} pour la gem et \texttt{target\_locations[2]} pour la clé.  

On ajoute également \texttt{is\_target\_locations\_visited} tableau de booléens de qui indique, pour chaque indice, si son $target\_tile$ associée à été visité ou pas.

\subsubsection{Modification de la boucle pour les \texttt{target\_tiles}}

Un fichier \texttt{target.c} a été ajouté pour gérer les \texttt{target\_tiles}.  

Dans la boucle de jeu, à chaque tour, les positions des \texttt{target\_tiles} sont triées par proximité au joueur à l'aide d'un tri à bulles sur leurs distances de Manhattan à travers \texttt{sort\_target\_locations}. Les emplacements déjà visités sont placés à la fin, et seules les positions non visitées sont affichées.  

De plus, lorsqu’un joueur se déplace sur un \texttt{target\_tile}, les ressources correspondantes lui sont attribuées grâce à \texttt{target\_location\_reached} qui vérifie si c’est le cas puis met à jour ses ressources.

\subsection{Ajouter une intelligence au joueur}

Un problème du mode non-itératif est que le joueur effectue des actions totalement aléatoires, conduisant souvent à des parties perdues.  

La solution est de créer un chemin qui approche le joueur de son objectif et le suivre.

Nous définissons \texttt{target\_location\_near} comme la position du \texttt{target\_tile} non visité le plus proche, ou la tuile de sortie si toutes ont été visitées. Un chemin vers cette position est alors calculé à l'aide de \texttt{path\_target\_finder}.

\subsubsection{path\_target\_finder}

La fonction \texttt{path\_target\_finder} prend en paramètres la localisation objective, la position actuelle du joueur contenue dans \texttt{game\_state}, ainsi qu’un buffer \texttt{path} qui est un tableau de localisations destiné à stocker le chemin calculé.

L’objectif de cette fonction est d’atteindre la localisation cible ou, à défaut, une tuile vide la plus proche possible de celle-ci, en empruntant un chemin de longueur minimale. Cette approche permet d’approcher progressivement le joueur de son objectif en minimisant le nombre de tuiles à poser. Pour cela, la fonction s’appuie sur une méthode récursive, \texttt{path\_target\_finder\_rec}, afin d’explorer les chemins possibles.

Avant de lancer la recherche récursive, les structures suivantes sont initialisées :
\begin{itemize}
  \item \texttt{dist}, représentant la distance de Manhattan entre la dernière tuile du chemin courant (souvent une tuile vide) et la tuile objectif ;
  \item \texttt{path\_len}, correspondant à la longueur du meilleur chemin trouvé ;
  \item \texttt{is\_visited\_position}, une matrice de la taille du manoir indiquant, pour chaque position, si elle a déjà été visitée afin d’éviter les cycles.
\end{itemize}

Ces variables sont ensuite passées par adresse à la fonction récursive et jouent le rôle de variables globales partagées entre tous les appels récursifs.

\subsubsection{Fonction récursive de recherche de chemin}

La fonction \texttt{path\_target\_finder\_rec} réalise une exploration en profondeur du plateau à partir de la position courante. À chaque appel, la position est ajoutée au chemin en cours et marquée comme visitée.\\
\\
\textbf{Cas de bas}:
Si la position courante correspond à la localisation objectif, le chemin courant est enregistré comme solution optimale. Si la longueur du chemin dépasse la taille maximale du plateau, l’exploration est interrompue. \\
\\
\textbf{Recursion}:
Pour chaque direction possible, la fonction vérifie que la position suivante est valide, non visitée et accessible selon les règles d’ouverture des tuiles. Deux cas sont alors distingués :
\begin{itemize}
  \item si la tuile suivante est vide, le chemin est évalué à l’aide de la distance de Manhattan à l’objectif et peut améliorer la meilleure solution courante
  \item si la tuile suivante est ouverte dans la direction opposée, la fonction s’appelle récursivement afin de poursuivre l’exploration
\end{itemize}

À la fin de l’exploration d’une position, celle-ci est dé-marquée, ce qui permet son utilisation dans d’autres chemins possibles.

\subsubsection{Complexité}

La fonction \texttt{path\_target\_finder} repose sur une exploration récursive en profondeur du plateau de jeu. Dans le pire des cas, chaque position du manoir est visitée au plus une fois grâce à la matrice \texttt{is\_visited\_position}, ce qui donne une complexité temporelle en \begin{scriptsize}O(BOARD\_SIZE)\end{scriptsize}.

  Cependant, en pratique, l’algorithme peut explorer plusieurs chemins partiels avant de converger vers une solution, ce qui rapproche son comportement d’un parcours exhaustif des chemins possibles. Le coût est donc exponentiel mais acceptable avec les contraintes du manoir.

La complexité spatiale est dominée par la matrice \texttt{is\_visited\_position} donc \begin{scriptsize}O(BOARD\_SIZE)\end{scriptsize}.

\subsubsection{Accommodation des changements dans la boucle de jeu}

Après le processus, nous avons un chemin \texttt{path}. Il amène vers la localisation de la tuile objective ou une tuile vide qui approche le joueur de l'objectif.

Dans la boucle de jeu en crée le chemin puis déplaçons le joueur au long du chemin jusqu'à atteindre \texttt{target\_location\_near} ou une tuile vide.

Si on a atteint \texttt{target\_location\_near} on recommence le processus

Sinon, on est devant une tuile vide. On place une tuile et recommençons le processus.

La limitation de cette fonction est que nous n'avons pas de contrôle sur le type de tuile placée. Il est possible que le joueur s'enferme en plaçant des terraces par exemple. S' il n'existe pas de chemins vers un objectif alors on fini le jeu et perdons.

La boucle est affichée dans les annexes.
